%!TEX program = lualatex
% SPDX-License-Identifier: CC-BY-SA-4.0
\documentclass[11pt]{article}

\usepackage{iftex}
\ifLuaTeX\else
  \PackageError{vintage-latex}{This example requires LuaLaTeX}
    {Run lualatex, not pdflatex or xelatex.}
\fi

\usepackage{fontspec}
\usepackage{xcolor}
\usepackage{geometry}
\usepackage[tracking=true,protrusion=false,expansion=false]{microtype}
\usepackage{tikz}

\defaultfontfeatures{Renderer=HarfBuzz}
\setmainfont{EB Garamond}[
  Numbers={OldStyle,Proportional},
  Ligatures={TeX,Common,Rare}
]

\definecolor{ragpaper}{HTML}{F3E8D5}
\definecolor{printerink}{HTML}{272019}
\pagecolor{ragpaper}
\color{printerink}
\geometry{a5paper,top=22mm,bottom=24mm,left=22mm,right=22mm}
\pagestyle{empty}
\setlength{\parindent}{1em}
\setlength{\emergencystretch}{.7em}

% A three-line initial built only from TeX primitives. The following word is
% a separate argument so it can be tracked small capitals.
\newcommand{\incipit}[2]{%
  \noindent
  \parshape 4
    3.15em \dimexpr\linewidth-3.15em\relax
    3.15em \dimexpr\linewidth-3.15em\relax
    3.15em \dimexpr\linewidth-3.15em\relax
    0pt \linewidth
  \makebox[0pt][r]{%
    \raisebox{-1.45\baselineskip}[0pt][0pt]{%
      \fboxsep=2.2pt\fboxrule=.45pt
      \fbox{\fontsize{39}{39}\selectfont #1}}%
    \hspace{.22em}}%
  \textls[85]{\textsc{#2}}%
}

\newcommand{\ornament}{%
  \begin{center}
  \begin{tikzpicture}[x=1cm,y=1cm,line width=.38pt,color=printerink]
    \draw (-2.25,0)--(-.45,0) ( .45,0)--(2.25,0);
    \draw (-.42,0) .. controls (-.22,.30) and (-.08,.30) .. (0,0)
                    .. controls ( .08,-.30) and ( .22,-.30) .. (.42,0);
    \draw (-.42,0) .. controls (-.22,-.30) and (-.08,-.30) .. (0,0)
                    .. controls ( .08,.30) and ( .22,.30) .. (.42,0);
    \fill (0,0) circle[radius=.65pt];
  \end{tikzpicture}
  \end{center}
}

\begin{document}

\begin{center}
  {\small\textls[135]{\textsc{A Short Account}}}\par
  \vspace{.35em}
  {\Large\textit{of the Air-Pump}}\par
  \vspace{.3em}
  {\footnotesize with observations made in a closed receiver}
\end{center}

\ornament

\incipit{W}{hen the receiver} had been exhausted by repeated strokes of the
piston, the enclosed bladder enlarged, although no hand was laid upon it. Upon
admitting the outward air, it returned by degrees to its former compass. This
alteration was entered in the journal before any cause was assigned.

The arrangement separates three effects that are often bundled together as
``old book style'': an ornament made from vector strokes, a boxed initial made
from ordinary TeX, and a narrow text measure. Each can be removed without
disturbing the other two.

\ornament

\begin{center}
  \footnotesize\textit{Experiment performed in the winter quarter, 1672.}
\end{center}

\vfill
\begin{center}
  {\scriptsize\textls[120]{\textsc{London: Printed for the Society}}}
\end{center}

\end{document}
